\section*{第 \ref{ch:g7:priorities} 章 --- 运算优先级}

\begin{solution}{exo:g7:priorities:1}
$5 + 3 \times 6 = 5 + 18 = 23$.

$(5 + 3) \times 6 = 8 \times 6 = 48$.

$30 - 12 \div 4 = 30 - 3 = 27$.

$(30 - 12) \div 4 = 18 \div 4 = 4.5$.
\end{solution}

\begin{solution}{exo:g7:priorities:2}
$24 - 6 + 2 = 18 + 2 = 20$（从左到右）。

$24 \div 6 \times 2 = 4 \times 2 = 8$（从左到右）。

$24 - 6 \times 2 = 24 - 12 = 12$（先算乘法）。

$24 \div (6 \times 2) = 24 \div 12 = 2$.
\end{solution}

\begin{solution}{exo:g7:priorities:3}
$7 + 3 \times (10 - 6) = 7 + 3 \times 4 = 7 + 12 = 19$.

$40 - (2 + 3) \times 4 = 40 - 5 \times 4 = 40 - 20 = 20$.

$60 \div \bigl(2 \times (7 - 4)\bigr) = 60 \div (2 \times 3)
= 60 \div 6 = 10$.
\end{solution}

\begin{solution}{exo:g7:priorities:4}
$\dfrac{18 + 6}{8} = \dfrac{24}{8} = 3$; \quad
$\dfrac{30}{7 - 2} = \dfrac{30}{5} = 6$; \quad
$\dfrac{5 \times 8}{4 + 6} = \dfrac{40}{10} = 4$.
\end{solution}

\begin{solution}{exo:g7:priorities:5}
$8 \times 102 = 8 \times (100 + 2) = 800 + 16 = 816$.

$25 \times 99 = 25 \times (100 - 1) = 2500 - 25 = 2475$.

$14 \times 12 + 14 \times 8 = 14 \times (12 + 8) = 14 \times 20 = 280$.

$37 \times 45 - 37 \times 35 = 37 \times (45 - 35) = 37 \times 10 =
370$.
\end{solution}

\begin{solution}{exo:g7:priorities:6}
“$50$ 与 $6$ 乘以 $7$ 之积的差”：
$50 - 6 \times 7 = 50 - 42 = 8$（不需要括号，优先级
自会起作用）。

“$4$ 与 $5$ 的和，与 $10$ 与
$8$ 的差之积”：$(4 + 5) \times (10 - 8) = 9 \times 2 = 18$。
\end{solution}

\begin{solution}{exo:g7:priorities:7}
$3 \times 2.50 + 2 \times 1.20 = 7.50 + 2.40 = 9.90$ 欧元。
\end{solution}

\begin{solution}{exo:g7:priorities:8}
(a) 不加括号：$4 + 2 \times 5 - 1 = 4 + 10 - 1 = 13$。

(b) $(4 + 2) \times 5 - 1 = 30 - 1 = 29$。

(c) $(4 + 2) \times (5 - 1) = 6 \times 4 = 24$。
\end{solution}

\begin{solution}{exo:g7:priorities:9}
$k = 1$：$12 + 4 \times 1 = 16$ 和 $16 \times 1 = 16$ --- 结果一致
（这正是 $k = 1$ 不能很好检验的原因！）$k = 2$：
$12 + 4 \times 2 = 20$，但 $16 \times 2 = 32$，两者不同。所以这两个
式子一般\emph{不}相等。错误的
等式 $12 + 4 \times k = 16 \times k$ 源于先算 $12 + 4$，
忽略了 $\times$ 的优先级。
\end{solution}

\begin{solution}{exo:g7:priorities:10}
\emph{1.} 团体价：$12 + 6 \times n$。当 $n = 5$ 时：
$12 + 30 = 42$ 欧元。

\emph{2.} 个人价：$9n$。比较：当 $n = 3$ 时，$30$ 与 $27$
（个人更便宜）；当 $n = 4$ 时，$36$ 与 $36$（相等）；当 $n = 5$ 时，$42$ 与
$45$（团体更便宜）。从 $5$ 人起，团体价更划算
（每增加一人，团体价多付 $6$ 而非 $9$，因此差距只会越来越大）。
\end{solution}

\begin{solution}{exo:g7:priorities:11}
可能的答案（答案不止一种）：
\[
24 = 1 \times 2 \times 3 \times 4, \qquad
10 = 1 + 2 + 3 + 4, \qquad
1 = (4 - 3) \times (2 - 1).
\]
\end{solution}

\begin{solution}{pb:g7:priorities:1}
\textbf{1.} $7 \times 102 = 7 \times (100 + 2) = 700 + 14 =
714$; \quad $9 \times 98 = 9 \times (100 - 2) = 900 - 18 = 882$.

\textbf{2.} $17 \times 13 + 17 \times 7 = 17 \times (13 + 7) =
17 \times 20 = 340$; \quad $45 \times 101 - 45 = 45 \times (101
- 1) = 45 \times 100 = 4\,500$.

\textbf{3.} $5 \times 87 = 870 \div 2 = 435$。合理性：
$5 = 10 \div 2$，因此乘以 $5$ 就是先乘以 $10$
再除以 $2$（$\times$ 与 $\div$ 属于同一级，
从左到右计算即可）。

\textbf{4.} $25 \times 32 = 3\,200 \div 4 = 800$，因为
$25 = 100 \div 4$。另有
$99 \times 47 = (100 - 1) \times 47 = 4\,700 - 47 = 4\,653$：
乘法对减法的分配律
（见\cref{thm:g7:priorities:distributivity}）。

\textbf{5.} $8 \times (5 + 3) = 64$ 与 $8 \times 5 + 8 \times 3
= 40 + 24 = 64$：两者相等 --- 这正是分配律。但
$8 \times (5 \times 3) = 8 \times 15 = 120$，而
$(8 \times 5) \times (8 \times 3) = 40 \times 24 = 960$：结果不相等。
因数分配的对象是\emph{和的各项}，而不是
积的各个因数。

\textbf{6.}
\begin{center}
\renewcommand{\arraystretch}{1.15}
\begin{tabular}{ccl}
$13$ & $24$ & 保留（$13$ 为奇数） \\
$6$ & $48$ & 划去（$6$ 为偶数） \\
$3$ & $96$ & 保留 \\
$1$ & $192$ & 保留 \\
\end{tabular}
\end{center}
保留下来的右列数之和：$24 + 96 + 192 = 312$，并且
确实有 $13 \times 24 = 312$。

\textbf{7.} 对 $21$ 逐次减半：$21, 10, 5, 2, 1$；对 $17$ 逐次加倍：
$17, 34, 68, 136, 272$。左列中的偶数 $10$ 与 $2$ 被划去；
保留的数之和为 $17 + 68 + 272 = 357$。验算：
$21 \times 17 = 357$。

\textbf{8.} 对 $16$ 逐次减半：$16, 8, 4, 2, 1$；对 $25$ 逐次加倍：
$25, 50, 100, 200, 400$。左列中除末尾的 $1$ 外，每个数都是
偶数，于是只剩一行保留，答案为 $400$。
若左列全由加倍得来，则只保留最后一行 ---
此时方法退化为单纯的加倍。

\textbf{9.} $13 = 12 + 1$，由分配律可得
$13 \times 24 = 12 \times 24 + 1 \times 24$；并且
$12 \times 24 = 6 \times 48$，因为一个因数减半而另一个因数加倍，
积保持不变
（$12 \times 24 = 6 \times 2 \times 24 = 6 \times 48$）。左列的奇数
$13$ 分出了右列数的一份：
第一行计入 $24$。

\textbf{10.} $6 \times 48 = 3 \times 96$（减半--加倍，
不计入：$6$ 为偶数）。$3 \times 96 = 2 \times 96 + 96 =
1 \times 192 + 96$：奇数 $3$ 计入 $96$。最后
$1 \times 192 = 192$：奇数 $1$ 计入 $192$，表格
结束。计入的数之和为 $24 + 96 + 192 = 312$。一行
计入其右列数，当且仅当其左列为奇数。因此，
“保留左列为奇数的行并求和”这一法则，恰是
分配律的记账过程。

\textbf{11.} $1 + 4 + 8 = 13$。对于 $21$，保留的行为
$17$、$68 = 4 \times 17$ 与 $272 = 16 \times 17$ 所对应的行：
即加倍数 $1 + 4 + 16 = 21$。

\textbf{12.} $25 = 16 + 8 + 1$ 与 $42 = 32 + 8 + 2$。

\textbf{13.} 取不超过
目标数的最大加倍数并减去之；余数更小，
且小于刚用过的加倍数（否则更大的加倍数本可放入）。
反复进行，余数严格变小，
最终必然到达 $0$：过程停止，所用的加倍数两两不同。
对于 $100$：$64$ 可放入（余 $36$），$32$ 可放入（余 $4$），
$4$ 可放入（余 $0$），因此 $100 = 64 + 32 + 4$。

\textbf{14.} $35$ 的加倍表：$1 \to 35$，$2 \to 70$，
$4 \to 140$，$8 \to 280$，$16 \to 560$。由于
$19 = 16 + 2 + 1$，取出这些行：
$560 + 70 + 35 = 665$。验算：$19 \times 35 = 665$。

\textbf{15.} 与书吏和农夫算法中的三个数相同：
$24 + 96 + 192$ --- 即 $24$ 的加倍数中，由
$13 = 1 + 4 + 8$ 选出的三项。把数写成加倍数之和，
正是用 $0$ 与 $1$ 书写它们；完整的讲述见
\cref{pb:g8:powers:1}。
\end{solution}
